\documentclass[conference,compsoc]{IEEEtran}

\usepackage[T1]{fontenc}
\usepackage[utf8]{inputenc}
\usepackage{cite}
\usepackage{graphicx}
\usepackage[caption=false,font=footnotesize]{subfig}
\usepackage{amsmath,amssymb,amsfonts,amsthm,mathtools,bm}
\usepackage{booktabs}
\usepackage{multirow}
\usepackage{enumitem}
\usepackage{xcolor}
\usepackage[hidelinks]{hyperref}
\usepackage{microtype}

\interdisplaylinepenalty=2500

\newtheorem{definition}{Definition}
\newtheorem{assumption}{Assumption}
\newtheorem{theorem}{Theorem}
\newtheorem{lemma}{Lemma}
\newtheorem{proposition}{Proposition}
\newtheorem{corollary}{Corollary}
\newtheorem{remark}{Remark}

\DeclareMathOperator{\Law}{Law}
\DeclareMathOperator{\clip}{clip}
\DeclareMathOperator{\supp}{supp}
\DeclareMathOperator{\Ber}{Bernoulli}
\DeclareMathOperator{\Bin}{Binomial}
\DeclareMathOperator{\argmax}{arg\,max}
\DeclareMathOperator{\argmin}{arg\,min}

\newcommand{\R}{\mathbb{R}}
\newcommand{\N}{\mathcal{N}}
\newcommand{\E}{\mathbb{E}}
\newcommand{\Prb}{\mathbb{P}}
\newcommand{\Ball}{\mathcal{B}}
\newcommand{\cZ}{\mathcal{Z}}
\newcommand{\cY}{\mathcal{Y}}
\newcommand{\cT}{\mathcal{T}}
\newcommand{\TV}{\mathrm{TV}}
\newcommand{\dd}{\,\mathrm d}
\newcommand{\norm}[1]{\left\lVert #1\right\rVert}
\newcommand{\inner}[2]{\left\langle #1,#2\right\rangle}

\title{Ball-Local Reconstruction Bounds for Nonconvex DP-SGD Transcripts}

\author{
Joseph Margaryan
}

\begin{document}
\maketitle

\begin{abstract}
We study reconstruction of one training record from a nonconvex DP-SGD
transcript under a local Ball prior on the hidden record.  The central
technical distinction is whether the adversary observes the Poisson inclusion
bits of the target record.  For the stronger known-inclusion transcript model,
we prove a Ball-local ReRo bound by dominating the adaptive transcript by a
product of revealed Bernoulli--Gaussian experiments.  For homogeneous training
schedules this product blow-up function is evaluated exactly as a finite
Binomial--Gaussian mixture.  For the weaker unknown-inclusion model, we define
the corresponding hidden-mixture f-DP/ReRo object and evaluate it by privacy
loss distribution convolution.  We keep this hidden f-DP curve distinct from
RDP-to-ReRo conversion bounds, which are certified but generally coarser and
apply to a different observation model than the revealed transcript.  Finally,
we give finite-prior exact-identification attacks whose likelihoods match the
known- and unknown-inclusion models.
\end{abstract}

\section{Setup}

Let
\[
D_z = D^- \cup \{z\},
\qquad
D_u = D^- \cup \{u\},
\]
where $D^-\in\cZ^{n-1}$ is fixed and $z$ is the hidden record drawn from a
prior $\pi_{u,r}$ supported on a radius-$r$ ball $\Ball(u,r)$.  A transcript
mechanism produces
\[
\tau = \cT_T(D)=(\tilde S_1,\ldots,\tilde S_T),
\]
where $\tilde S_t$ is the sanitized DP-SGD update at step $t$.  A reconstructor
is a measurable map
\[
R:\Theta\times \cZ^{n-1}\to \cZ.
\]
For a loss $\rho$ and tolerance $\eta$, the reconstruction success probability is
\[
p_{\rm succ}(\eta)
=
\Pr_{Z\sim\pi_{u,r},\,\tau\sim \Law(\cT_T(D_Z))}
\left[
\rho(Z,R(\tau,D^-))\le \eta
\right].
\]

\begin{definition}[Prior anti-concentration]
For a prior $\pi$ and loss $\rho$, define
\[
\kappa_{\pi,\rho}(\eta)
:=
\sup_{z_0\in\cZ}
\Pr_{Z\sim\pi}\left[\rho(Z,z_0)\le \eta\right].
\]
For a uniform prior on a Euclidean $d$-ball of radius $r$ and Euclidean loss,
ignoring boundary effects at the maximizing center,
\[
\kappa(\eta)=
\begin{cases}
(\eta/r)^d, & 0\le \eta<r,\\
1, & \eta\ge r.
\end{cases}
\]
\end{definition}

\section{Hypothesis-testing notation}

\begin{definition}[Blow-up function]
For probability laws $P,Q$ on a common measurable space,
\[
B_\kappa(P\|Q)
:=
\sup_{A:\,Q(A)\le \kappa} P(A),
\qquad
\kappa\in[0,1].
\]
Equivalently, if $T_{Q,P}$ is the tradeoff function with type-I error measured
under $Q$ and power under $P$, then
\[
B_\kappa(P\|Q)=1-T_{Q,P}(\kappa).
\]
\end{definition}

The order matters.  Throughout this paper, $Q$ is the reference/null law and
$P$ is the candidate/alternative law.  Thus a reconstruction event whose
reference probability is at most $\kappa$ has candidate probability at most
$B_\kappa(P\|Q)$.

\begin{lemma}[Post-processing]
If $K$ is a probability kernel, then
\[
B_\kappa(PK\|QK)\le B_\kappa(P\|Q),
\qquad
\kappa\in[0,1].
\]
\end{lemma}

\begin{proof}
Every event $A$ in the output space induces the randomized test
$x\mapsto K(x,A)$ in the input space.  If $(QK)(A)\le \kappa$, then this test
has $Q$-mass at most $\kappa$, and its $P$-mass is $(PK)(A)$.  Taking the
supremum over output events gives the claim.
\end{proof}

\section{A local Ball-ReRo theorem}

For each $z\in\supp(\pi_{u,r})$ write
\[
P_z := \Law(\cT_T(D_z)),
\qquad
Q_u := \Law(\cT_T(D_u)).
\]

\begin{theorem}[Local Ball-ReRo]
Let $\pi=\pi_{u,r}$ and suppose there is a pair $(P^\star,Q^\star)$ such that
for all $z\in\supp(\pi)$ and all $\kappa\in[0,1]$,
\[
B_\kappa(P_z\|Q_u)
\le
g(\kappa)
:=
B_\kappa(P^\star\|Q^\star),
\]
where $g$ is nondecreasing and concave.  Then every measurable reconstructor
satisfies
\[
p_{\rm succ}(\eta)
\le
g\!\left(\kappa_{\pi,\rho}(\eta)\right).
\]
\end{theorem}

\begin{proof}
For fixed $z$, define
\[
E_z := \{\tau:\rho(z,R(\tau,D^-))\le\eta\}
\]
and
\[
\kappa_z:=Q_u(E_z).
\]
Then
\[
P_z(E_z)
\le
B_{\kappa_z}(P_z\|Q_u)
\le
g(\kappa_z).
\]
Averaging over $Z\sim\pi$ gives
\[
p_{\rm succ}(\eta)
\le
\E_{Z\sim\pi}g(\kappa_Z).
\]
By concavity,
\[
\E g(\kappa_Z)
\le
g(\E\kappa_Z).
\]
It remains to bound $\E\kappa_Z$.  Let
\[
A=\{(z,\tau):\rho(z,R(\tau,D^-))\le\eta\}.
\]
By Tonelli's theorem,
\[
\E_Z\kappa_Z
=
\int Q_u(E_z)\,\pi(\dd z)
=
\int
\Pr_{Z\sim\pi}\!\left[\rho(Z,R(\tau,D^-))\le\eta\right]
Q_u(\dd\tau).
\]
For fixed $\tau$, $R(\tau,D^-)$ is a deterministic point of $\cZ$, so the inner
probability is at most $\kappa_{\pi,\rho}(\eta)$.  Hence
\[
\E_Z\kappa_Z \le \kappa_{\pi,\rho}(\eta).
\]
Since $g$ is nondecreasing,
\[
p_{\rm succ}(\eta)
\le
g(\kappa_{\pi,\rho}(\eta)).
\]
\end{proof}

\section{Ball-SGD transcript domination}

At step $t$, conditional on the past history $h$, the sanitized update can be
written as
\[
\tilde S_t(D_u)
=
C_t^h + I_t \bar g_t^h(u) + \xi_t,
\]
\[
\tilde S_t(D_z)
=
C_t^h + I_t \bar g_t^h(z) + \xi_t,
\]
where $I_t\sim\Ber(q_t)$ is the Poisson inclusion bit and
$\xi_t\sim\N(0,\sigma_t^2 I_p)$.  Define the clipped gradient gap
\[
\delta_t^h(z,u)
:=
\bar g_t^h(z)-\bar g_t^h(u),
\]
and the normalized Ball sensitivity
\[
c_t
:=
\frac{\Delta_t(r)}{\sigma_t},
\qquad
\Delta_t(r)
\ge
\sup_{z\in\supp(\pi_{u,r}),h}
\norm{\delta_t^h(z,u)}_2.
\]
The Ball-specific improvement is that $\Delta_t(r)$ may be much smaller than
the standard clipping-only sensitivity $2C_t$.

\subsection{Known inclusion: revealed transcript bound}

The known-inclusion adversary observes the target inclusion bits or an
equivalent residualized transcript from which the selected target steps are
known.  The one-step revealed reference pair is
\[
Q_t^{\rm rev}
=
\Law(I_t,G_t),
\qquad
P_t^{\rm rev}
=
\Law(I_t,I_t c_t+G_t),
\]
where $G_t\sim\N(0,1)$ and $I_t\sim\Ber(q_t)$ are independent.  Define
\[
Q_{1:T}^{\rm rev}:=\bigotimes_{t=1}^T Q_t^{\rm rev},
\qquad
P_{1:T}^{\rm rev}:=\bigotimes_{t=1}^T P_t^{\rm rev}.
\]

\begin{theorem}[Known-inclusion Ball-ReRo]
For the revealed-inclusion transcript model,
\[
p_{\rm succ}(\eta)
\le
B_{\kappa_{\pi,\rho}(\eta)}
\!\left(P_{1:T}^{\rm rev}\middle\|Q_{1:T}^{\rm rev}\right).
\]
\end{theorem}

\begin{proof}
For fixed $z$ and history $h$, center by the common non-target contribution and
project onto the direction of $\delta_t^h(z,u)$.  Conditional on $I_t=i$, the
two laws differ by a one-dimensional Gaussian shift of magnitude
$i\|\delta_t^h(z,u)\|_2/\sigma_t\le i c_t$.  A Gaussian contraction kernel maps
the worst-case shift $i c_t$ to the actual shift.  Adding the common center and
orthogonal Gaussian coordinates is another kernel.  These one-step kernels glue
adaptively over $t$, giving a global kernel $K_z$ such that
\[
P_{1:T}^{\rm rev}K_z=P_z,
\qquad
Q_{1:T}^{\rm rev}K_z=Q_u.
\]
By post-processing,
\[
B_\kappa(P_z\|Q_u)
\le
B_\kappa(P_{1:T}^{\rm rev}\|Q_{1:T}^{\rm rev}).
\]
The local Ball-ReRo theorem then gives the result.
\end{proof}

\subsection{Exact homogeneous evaluator}

When $q_t=q$ and $c_t=c$ for all $t$, the privacy loss
\[
L=\log\frac{\dd P_{1:T}^{\rm rev}}{\dd Q_{1:T}^{\rm rev}}
\]
has a finite Binomial--Gaussian mixture form.  Let
\[
K:=\sum_{t=1}^T I_t\sim\Bin(T,q).
\]
Conditional on $K=k$,
\[
L\mid K=k,Q
\sim
\N\!\left(-\frac12 k c^2,\,k c^2\right),
\]
and
\[
L\mid K=k,P
\sim
\N\!\left(+\frac12 k c^2,\,k c^2\right),
\]
with an atom at $L=0$ when $k=0$.  The Neyman--Pearson event is an upper level
set of $L$.  Therefore
\[
B_\kappa(P_{1:T}^{\rm rev}\|Q_{1:T}^{\rm rev})
=
P[L>\tau_\kappa]+\theta P[L=\tau_\kappa],
\]
where $\tau_\kappa$ and $\theta\in[0,1]$ are chosen so that
\[
Q[L>\tau_\kappa]+\theta Q[L=\tau_\kappa]=\kappa.
\]
This is the exact evaluator used for the known-inclusion curve in the
experiments.

\subsection{Unknown inclusion: hidden-mixture f-DP bound}

If the transcript does not reveal target inclusion bits, the one-step hidden
reference pair is
\[
Q_t^{\rm hid}=\N(0,1),
\qquad
P_t^{\rm hid}=(1-q_t)\N(0,1)+q_t\N(c_t,1).
\]
The product hidden f-DP object is
\[
B_\kappa(P_{1:T}^{\rm hid}\|Q_{1:T}^{\rm hid}),
\qquad
P_{1:T}^{\rm hid}=\bigotimes_t P_t^{\rm hid},
\quad
Q_{1:T}^{\rm hid}=\bigotimes_t Q_t^{\rm hid}.
\]
For one step, the privacy loss under $Q_t^{\rm hid}$ is
\[
\ell_t(x)
=
\log\left((1-q_t)+q_t\exp(c_t x-c_t^2/2)\right),
\qquad
x\sim\N(0,1).
\]
The product privacy loss is $\sum_t \ell_t(X_t)$.  We evaluate its distribution
by Gauss--Hermite quadrature and convolution.  This gives a deterministic
hypothesis-testing approximation to the hidden-inclusion f-DP curve.  It is not
mixed with the revealed-inclusion bound: the observation models differ.

\section{RDP-to-ReRo conversion}

If a mechanism satisfies an RDP bound
\[
D_\alpha(P\|Q)\le \varepsilon_\alpha,
\qquad \alpha>1,
\]
then H\"older's inequality gives
\[
B_\kappa(P\|Q)
\le
\min\left\{
1,\,
\exp\left(
\frac{\alpha-1}{\alpha}
\left[\log\kappa+\varepsilon_\alpha\right]
\right)
\right\}.
\]
Optimizing over orders,
\[
\Gamma_{\rm RDP}(\kappa)
=
\inf_{\alpha>1}
\min\left\{
1,\,
\exp\left(
\frac{\alpha-1}{\alpha}
\left[\log\kappa+\varepsilon_\alpha\right]
\right)
\right\}.
\]
This conversion is a certified bound for the law pair accounted for by the RDP
accountant.  In our experiments it is reported as a hidden-subsampling/final
release certificate.  It should not be interpreted as a bound on the stronger
revealed-inclusion transcript unless the RDP accountant explicitly includes the
released inclusion bits.

\section{Finite-prior exact-identification attacks}

Let $\{z_1,\ldots,z_m\}$ be a public candidate support with prior weights
$\pi_i$.  The exact-identification baseline is
\[
\kappa_{\rm id}=\max_i \pi_i,
\]
and for a uniform finite prior, $\kappa_{\rm id}=1/m$.

After residualizing the known non-target batch contribution, a known-inclusion
Gaussian transcript has observations
\[
r_t = g_t(z_i)+\xi_t
\]
on selected target-present steps.  The exact finite-prior Bayes score is
\[
S_i^{\rm known}
=
\log \pi_i
-
\frac12
\sum_{t\in\mathcal I}
\frac{\norm{r_t-g_t(z_i)}_2^2}{\sigma_t^2},
\]
with the obvious least-squares limit when $\sigma_t=0$.

For unknown inclusion,
\[
r_t =
\begin{cases}
\xi_t, & I_t=0,\\
g_t(z_i)+\xi_t, & I_t=1,
\end{cases}
\]
so the exact score is
\[
S_i^{\rm unknown}
=
\log\pi_i
+
\sum_t
\log\left[
(1-q_t)\exp\left(-\frac{\norm{r_t}_2^2}{2\sigma_t^2}\right)
+
q_t\exp\left(-\frac{\norm{r_t-g_t(z_i)}_2^2}{2\sigma_t^2}\right)
\right],
\]
up to candidate-independent constants.  The attack returns
\[
\hat i=\argmax_i S_i.
\]
The known-inclusion attack should be compared against the revealed f-DP curve;
the unknown-inclusion attack should be compared against the hidden f-DP/RDP
curves.

\section{Relationship to prior ReRo work}

Hayes, Mahloujifar, and Balle formulate reconstruction risk through the blow-up
function and estimate DP-SGD reconstruction bounds by sampling or otherwise
approximating the relevant likelihood-ratio distribution.  Our known-inclusion
contribution specializes this philosophy to a Ball-local transcript pair and
gives an exact homogeneous evaluator for the revealed Bernoulli--Gaussian
product reference experiment.

Kaissis, Hayes, Ziller, and Rueckert emphasize the hypothesis-testing
interpretation of reconstruction bounds and the distinction between analytic
mechanism curves and numerical approximations.  Our implementation follows this
separation: the exact homogeneous revealed curve is a theorem-backed numerical
evaluation, the hidden f-DP curve is a deterministic PLD approximation of a
different hypothesis-testing object, and the RDP conversion is a separate
certified conversion bound.

\section{Experimental protocol}

The synthetic experiment trains three models on linearly separable binary data:
nonprivate ERM, Ball-DP, and standard DP calibrated to the same target privacy
level.  The main figures are:
\begin{enumerate}[leftmargin=*]
\item utility over training and final utility;
\item known-inclusion revealed f-DP curves for Ball-DP and standard DP;
\item unknown-inclusion hidden f-DP and RDP conversion curves;
\item a monotone prefix privacy--utility trajectory for the known-inclusion
      transcript bound;
\item finite-prior exact-identification bars for known- and unknown-inclusion
      attacks.
\end{enumerate}
The direct stepwise composition bound is valid but vacuous in the default
high-$T$ experiment, so it is kept as a diagnostic rather than a main figure.

\section{Conclusion}

The main methodological lesson is that reconstruction certificates must be
indexed by the observation model.  A revealed-inclusion transcript bound, a
hidden-inclusion f-DP curve, and an RDP conversion bound are all useful, but they
are not interchangeable.  Ball-local sensitivity gives meaningful
reconstruction certificates in the revealed transcript model where standard
clipping-based sensitivity is nearly vacuous, and it also supports exact
finite-prior attacks whose likelihoods match the corresponding threat model.

\bibliographystyle{IEEEtran}
\begin{thebibliography}{9}

\bibitem{abadi2016}
M.~Abadi et al.,
``Deep Learning with Differential Privacy,''
in \emph{Proc. ACM CCS}, 2016.

\bibitem{mironov2017}
I.~Mironov,
``R\'enyi Differential Privacy,''
in \emph{Proc. IEEE CSF}, 2017.

\bibitem{dong2022}
J.~Dong, A.~Roth, and W.~Su,
``Gaussian Differential Privacy,''
\emph{Journal of the Royal Statistical Society: Series B}, 2022.

\bibitem{balle2022}
B.~Balle, G.~Cherubin, and J.~Hayes,
``Reconstructing Training Data with Informed Adversaries,''
in \emph{Proc. IEEE S\&P}, 2022.

\bibitem{hayes2023}
J.~Hayes, S.~Mahloujifar, and B.~Balle,
``Bounding Training Data Reconstruction in DP-SGD,''
in \emph{Advances in Neural Information Processing Systems}, 2023.

\bibitem{kaissis2023}
G.~Kaissis, J.~Hayes, A.~Ziller, and D.~Rueckert,
``Bounding Data Reconstruction Attacks with the Hypothesis Testing
Interpretation of Differential Privacy,''
arXiv:2307.03928, 2023.

\end{thebibliography}

\end{document}
